\section*{Chapter \ref{ch:b1:enzymes} --- Enzymes and Biochemical Catalysis}

\begin{solution}{exo:b1:enzymes:1}
It lowers the \omterm{prop:b1:enzymes:whatitdoes}{activation energy} (the height of the \omterm{prop:b1:enzymes:whatitdoes}{transition state}),
and so the rate, in both directions; it \omterm{def:b1:flowering-plant-organization:organs}{leaves} $\Delta G$, the
equilibrium constant and the position of equilibrium unchanged. On the
profile the peak is lowered, the two ends are not.
\end{solution}

\begin{solution}{exo:b1:enzymes:2}
$K_m$: \omterm{def:b1:enzymes:enzyme}{substrate} concentration at half $V_{\max}$ (\unit{mol/L}).
$V_{\max}$: rate at saturating \omterm{def:b1:enzymes:enzyme}{substrate} (\unit{mol/s}, or \unit{\micro
mol/min}). $k_{\text{cat}}$: \omterm{def:b1:enzymes:enzyme}{substrate} molecules converted per second per
\omterm{def:b1:enzymes:enzyme}{enzyme} at saturation (\unit{s^{-1}}). $k_{\text{cat}}/K_m$: efficiency,
the second-order rate constant at low \omterm{def:b1:enzymes:enzyme}{substrate}
(\unit{L.mol^{-1}.s^{-1}}).
\end{solution}

\begin{solution}{exo:b1:enzymes:3}
No \omterm{def:b1:enzymes:inhibitors}{inhibitor}: intercept 1, so $V_{\max} = 1$; $x$ intercept $-0.5$, so
$K_m = 2$. Competitive: $V_{\max} = 1$, $x$ intercept $-0.25$, $K_m = 4$.
Non-competitive: intercept 2, $V_{\max} = 0.5$; $x$ intercept $-0.5$,
$K_m = 2$.
\end{solution}

\begin{solution}{exo:b1:enzymes:4}
\omterm{def:b1:proteins:allostery}{Allosteric} (feedback) regulation, milliseconds; \omterm{prop:b1:enzymes:control}{covalent modification}
(phosphorylation), seconds to minutes; proteolytic activation of a
zymogen, once, on demand; control of the amount by gene expression and
degradation, minutes to hours.
\end{solution}

\begin{solution}{exo:b1:enzymes:5}
$v_0 = 50[S]/(0.1 + [S])$: 8.3, 25, 45.5 and \qty{49.5}{\micro mol/min}.
$0.9 = [S]/(0.1 + [S])$ gives $[S] = 0.9\,K_m/0.1 = \qty{0.9}{mmol/L}$.
\end{solution}

\begin{solution}{exo:b1:enzymes:6}
$e^{20\,000/2580} = e^{7.75} = 2300$. For $10^{10}$: $\Delta E_a =
RT\ln 10^{10} = 2.58\times 23.0 = \qty{59}{kJ/mol}$.
\end{solution}

\begin{solution}{exo:b1:enzymes:7}
$V_{\max} = \qty{0.12}{mol/L}$ per minute in \qty{1}{mL}:
\qty{1.2e-4}{mol/min} $= \qty{2.0e-6}{mol/s}$ of product; \omterm{def:b1:enzymes:enzyme}{enzyme}
\qty{2e-12}{mol}: $k_{\text{cat}} = 2\times 10^{-6}/2\times 10^{-12} =
\qty{1e6}{s^{-1}}$. Efficiency $10^6/0.012 = \qty{8e7}{L.mol^{-1}.s^{-1}}$,
within a factor of ten of the diffusion limit: nearly every encounter
with a \omterm{def:b1:enzymes:enzyme}{substrate} molecule is productive.
\end{solution}

\begin{solution}{exo:b1:enzymes:8}
$1 + 2/K_i = 2$: $K_i = \qty{2}{mmol/L}$. For \qty{90}{\%}:
$[S] = 9K_m^{\text{app}}$: \qty{9}{mmol/L} without, \qty{18}{mmol/L}
with the \omterm{def:b1:enzymes:inhibitors}{inhibitor}.
\end{solution}

\begin{solution}{exo:b1:enzymes:9}
Inhibiting the first committed step stops the whole pathway at once
and wastes no intermediates, which would otherwise accumulate; the
end product is the signal that matters --- its abundance is exactly
the information that the pathway is no longer needed --- whereas an
intermediate's level says nothing about demand.
\end{solution}

\begin{solution}{exo:b1:enzymes:10}
A protease active in the \omterm{def:b1:cell-unit-of-life:cell}{cell} that makes it would digest that \omterm{def:b1:cell-unit-of-life:cell}{cell};
made as an inactive zymogen, it is harmless until enteropeptidase, an
\omterm{def:b1:enzymes:enzyme}{enzyme} of the intestinal lining, cuts trypsinogen to trypsin in the
duodenum, where digestion is wanted; trypsin then activates the other
zymogens (a cascade). Because a trace of trypsin can start the cascade
prematurely, the pancreas packs a specific trypsin \omterm{def:b1:enzymes:inhibitors}{inhibitor} into its
granules as insurance; failure of the arrangement is pancreatitis.
\end{solution}

\begin{solution}{exo:b1:enzymes:11}
Without: $3^3/(3^3 + 3^3) = 0.50$. With: $27/(166 + 27) = 0.14$: the
rate falls to \qty{28}{\%} of its value. A Michaelis--Menten \omterm{def:b1:enzymes:enzyme}{enzyme}
with $K_m$ raised from 3 to 5.5: $3/6 = 0.50$ to $3/8.5 = 0.35$, i.e.\
\qty{70}{\%}. \omterm{def:b1:proteins:allostery}{Cooperativity} makes the same shift of the curve more
than twice as effective near the working point: a switch rather than
a dimmer.
\end{solution}

\begin{solution}{exo:b1:enzymes:12}
A bare catalyst (a platinum surface, an acid) speeds many reactions
indiscriminately; an \omterm{def:b1:enzymes:enzyme}{enzyme} speeds one reaction on one \omterm{def:b1:enzymes:enzyme}{substrate} ---
the specificity of its \omterm{def:b1:enzymes:enzyme}{active site} is already information about what
the \omterm{def:b1:cell-unit-of-life:cell}{cell} wants done. Regulation adds a second layer: \omterm{def:b1:proteins:allostery}{allosteric} sites,
phosphorylation sites and zymogen peptides tell the \omterm{def:b1:enzymes:enzyme}{enzyme} when and
where to act, in response to the \omterm{def:b1:cell-unit-of-life:cell}{cell}'s state. The cost is a large
\omterm{def:b1:proteins:peptide}{protein} (hundreds of residues for a site of a dozen), slower turnover
than the best inorganic catalysts, and the machinery to make, modify
and destroy it; the benefit is a chemistry that runs only where and
when it is useful, which is what a metabolism is.
\end{solution}

\begin{solution}{pb:b1:enzymes:1}
\textbf{1.} $1/[S]$: 2, 1, 0.5, 0.2, 0.1, 0.05. $1/v_0$: 5.0, 3.0, 2.0,
1.4, 1.2, 1.1.
\textbf{2.} Each step of $1/[S]$ changes $1/v_0$ in proportion: slope
$(5.0 - 1.1)/(2 - 0.05) = 2.0$; intercept $1.0$.
\textbf{3.} $V_{\max} = 1/1.0 = \qty{1.0}{\micro mol/min}$; $K_m =
\text{slope}\times V_{\max} = \qty{2.0}{mmol/L}$.
\textbf{4.} $1.0\times 5/(2 + 5) = 0.714$: as measured.
\textbf{5.} $\qty{10e-9}{mol/L}\times\qty{1e-3}{L} = \qty{1e-11}{mol}$;
$V_{\max} = \qty{1e-6}{mol/min} = \qty{1.67e-8}{mol/s}$;
$k_{\text{cat}} = 1.67\times 10^{-8}/10^{-11} = \qty{1670}{s^{-1}}$.
\textbf{6.} $1670/(2\times 10^{-3}) = \qty{8.3e5}{L.mol^{-1}.s^{-1}}$:
a thousand times below the diffusion limit; only one collision in a
thousand is productive.
\textbf{7.} $0.95 = [S]/(2 + [S])$: $[S] = 19\,K_m = \qty{38}{mmol/L}$.
\textbf{8.} $1/v_0$: 9.0, 5.0, 3.0, 1.8, 1.4, 1.2. Slope $(9.0 -
1.2)/1.95 = 4.0$; intercept $1.0$.
\textbf{9.} $V_{\max}^{\text{app}} = \qty{1.0}{\micro mol/min}$
(unchanged); $K_m^{\text{app}} = 4.0\times 1.0 = \qty{4.0}{mmol/L}$.
\textbf{10.} $K_m$ doubled, $V_{\max}$ unchanged: competitive.
\textbf{11.} $4 = 2(1 + 1/K_i)$: $K_i = \qty{1.0}{mmol/L}$.
\textbf{12.} With the \omterm{def:b1:enzymes:inhibitors}{inhibitor}, $v = V_{\max}[S]/(K_m(1 + [I]/K_i) +
[S])$; halving the rate requires $K_m(1 + [I]/K_i) + [S] = 2(K_m +
[S])$, i.e.\ $[I] = K_i(K_m + [S])/K_m$: at \qty{2}{mmol/L}, $[I] =
1\times 4/2 = \qty{2}{mmol/L}$; at \qty{20}{mmol/L}, $[I] = 22/2 =
\qty{11}{mmol/L}$. The more \omterm{def:b1:enzymes:enzyme}{substrate}, the more \omterm{def:b1:enzymes:inhibitors}{inhibitor} it takes.
\textbf{13.} A molecule resembling the \omterm{def:b1:enzymes:enzyme}{substrate} (or its \omterm{prop:b1:enzymes:whatitdoes}{transition
state}) --- an ester analogue that cannot be hydrolysed --- binding in
the \omterm{def:b1:enzymes:enzyme}{active site}.
\textbf{14.} $1/v_0$: 10.0, 6.0, 4.0, 2.8, 2.4, 2.2; slope $4.0$,
intercept $2.0$: $V_{\max}^{\text{app}} = \qty{0.5}{\micro mol/min}$,
$K_m^{\text{app}} = 4.0\times 0.5 = \qty{2.0}{mmol/L}$.
\textbf{15.} $V_{\max}$ halved, $K_m$ unchanged: non-competitive;
$2 = 1 + 1/K_i$: $K_i = \qty{1.0}{mmol/L}$.
\textbf{16.} J binds at a site other than the \omterm{def:b1:enzymes:enzyme}{active site}, on both
free \omterm{def:b1:enzymes:enzyme}{enzyme} and $ES$, with the same affinity: \omterm{def:b1:enzymes:enzyme}{substrate} does not
compete with it, and at any $[S]$ half the \omterm{def:b1:enzymes:enzyme}{enzyme} molecules are
inactivated.
\textbf{17.} With J the rate doubles (half of twice as much \omterm{def:b1:enzymes:enzyme}{enzyme} is
still active); with I it also doubles (the rate is proportional to
\omterm{def:b1:enzymes:enzyme}{enzyme} at every $[S]$): doubling the \omterm{def:b1:enzymes:enzyme}{enzyme} never tells the \omterm{def:b1:enzymes:inhibitors}{inhibitors}
apart.
\textbf{18.} Rate at \qty{25}{\celsius}: $V_{\max}[S]/(K_m + [S])$ with
$V_{\max} = k_{\text{cat}}[E]$: per litre, $1670\times 10^{-8} =
\qty{1.67e-5}{mol/L/s}$; $\times 0.5/2.5 = \qty{3.3e-6}{mol/L/s}$; at
\qty{37}{\celsius}, $\times 2^{1.2} = 2.3$: \qty{7.7e-6}{mol/L/s}; in
\qty{1e-12}{L}: \qty{7.7e-18}{mol/s} $= \num{4.6e6}$ molecules per
second.
\textbf{19.} No rise needed: the rate (\num{4.6e6}) exceeds the demand
(\num{3e6}); if it had to rise, the \omterm{def:b1:cell-unit-of-life:cell}{cell} could raise the \omterm{def:b1:enzymes:enzyme}{substrate}, or
phosphorylate the \omterm{def:b1:enzymes:enzyme}{enzyme} to lower its $K_m$, or make more \omterm{def:b1:enzymes:enzyme}{enzyme}.
\textbf{20.} $K_m^{\text{app}} = 2(1 + 0.4/0.2) = \qty{6}{mmol/L}$: rate
$\times 0.5/6.5$ instead of $0.5/2.5$: \qty{38}{\%} of before,
\num{1.8e6} per second --- the product throttles its own synthesis to
below the demand, so it is consumed and its level falls, which releases
the \omterm{def:b1:enzymes:enzyme}{enzyme}: a self-adjusting supply.
\textbf{21.} Before: $0.5/2.5 = 0.20$ of $V_{\max}$; after: $0.5/0.9 =
0.56$: nearly threefold.
\textbf{22.} Near zero: at pH 5 the histidine of the \omterm{def:b1:enzymes:enzyme}{active site} is
protonated and cannot act as a base, and acidic residues of the site
are neutralised; the \omterm{def:b1:enzymes:enzyme}{enzyme}'s ionisation state, and perhaps its fold,
are wrong. It is built for the \omterm{def:b1:eukaryotic-cell:organelle}{cytosol}, not the \omterm{def:b1:eukaryotic-cell:endomembrane}{lysosome}.
\textbf{23.} $K_m$ changes little (binding is mostly by the rest of the
site); $k_{\text{cat}}$ collapses by orders of magnitude, since the
histidine was the catalytic base that activated the water or serine.
\textbf{24.} Near $K_m$ the rate responds to changes of \omterm{def:b1:enzymes:enzyme}{substrate}
(slope near $V_{\max}/2K_m$); far above it the \omterm{def:b1:enzymes:enzyme}{enzyme} is saturated and
insensitive, and the excess \omterm{def:b1:enzymes:enzyme}{substrate} would be an osmotic and
chemical burden. Working near $K_m$ makes the pathway controllable by
supply.
\textbf{25.} $K_m = \qty{2.0}{mmol/L}$, $V_{\max} = \qty{1.0}{\micro mol/min}$
(\qty{10}{nmol/L} \omterm{def:b1:enzymes:enzyme}{enzyme}), $k_{\text{cat}} = \qty{1670}{s^{-1}}$; I
competitive, $K_i = \qty{1.0}{mmol/L}$; J non-competitive, $K_i =
\qty{1.0}{mmol/L}$.
\end{solution}
